\section{Empirical Comparison}
\label{sec:comparison}

\subsection{Setup}

We compare three search strategies:
\begin{enumerate}
  \item \textbf{Single-scale \recom{}} (\texttt{--search single}):
    standard \recom{} with no coarse level, single seed.
  \item \textbf{Fixed-$\alpha$ \ms{}} (\texttt{--search multiscale
    --multiscale-alpha 0.30}): multi-scale chain with $\alpha = 0.30$
    for all states.
  \item \textbf{\ams{}} (\texttt{--search multiscale-adaptive}): adaptive
    chain with $\alpha_0 = 0.30$, $\tau = 0.30$, $\Delta = 50$,
    $\gamma_0 = 0.10$.
\end{enumerate}

All runs use $T = 5{,}000$ steps, the geographic edge-weight signal
(\texttt{--weights geographic}), and a prime-factor structure
(\texttt{--structure prime-factor}).
We measure:
\begin{itemize}
  \item \textbf{ACF@100}: autocorrelation of the $\EC$ sequence at lag 100.
    Lower ACF indicates faster mixing.
  \item \textbf{Final alpha} (for \ams{} only): the converged $\alpha^*$.
  \item \textbf{Runtime}: wall-clock seconds on a single core.
\end{itemize}

\subsection{Results}

\begin{table}[h]
  \centering
  \caption{Comparison of search strategies by state. ACF@100 is the
    autocorrelation of the edge-cut objective at lag 100; lower is better.
    Runtime is single-core wall time for $T = 5{,}000$ steps.}
  \label{tab:comparison}
  \begin{tabular}{llrrrrr}
    \toprule
    State & $k$ & \multicolumn{2}{c}{ACF@100} & & $\alpha^*$ & Runtime \\
    \cmidrule{3-4}
    & & \recom{} & MS fixed & MS adaptive & (AMS) & (AMS, s) \\
    \midrule
    NC & 14 & 0.71 & 0.52 & 0.50 & 0.28 & 410 \\
    TX & 38 & 0.83 & 0.61 & 0.58 & 0.33 & 1{,}820 \\
    IA &  4 & 0.44 & 0.51 & 0.41 & 0.17 & 95 \\
    WI &  8 & 0.62 & 0.55 & 0.53 & 0.26 & 280 \\
    CA & 52 & 0.89 & 0.74 & 0.69 & 0.22 & 3{,}450 \\
    \bottomrule
  \end{tabular}
\end{table}

\subsection{Interpretation}

\paragraph{NC, TX, WI: marginal improvement.}
For states where $\alpha = 0.30$ is near-optimal, \ams{} converges to a
similar $\alpha^*$ and achieves comparable ACF to fixed-$\alpha$ \ms{}.
The improvement over \recom{} is similar in both \ms{} variants, confirming
that the coarse level (not the alpha tuning) drives the mixing gain.
The adaptive overhead --- the extra bookkeeping of the window and alpha update
--- is negligible compared to the chain step cost.

\paragraph{IA: clear improvement from adaptation.}
Iowa is the most instructive case.
Fixed-$\alpha$ \ms{} with $\alpha = 0.30$ \emph{worsens} mixing compared to
single-scale \recom{} (ACF 0.51 vs.\ 0.44): too many coarse proposals fail
the rebalancer, injecting rejections that stall the chain.
\ams{} detects this via the low coarse acceptance rate, drives $\alpha$ down
to 0.17, and achieves ACF 0.41 --- better than both single-scale and fixed
\ms{}.
This is the central case for the adaptive approach: fixed-$\alpha$ \ms{}
can hurt if $\alpha$ is miscalibrated.

\paragraph{CA: largest improvement.}
California has the largest $k$ and the most extreme county population
heterogeneity (Los Angeles County alone holds $\approx 25\%$ of the state
population).
The rebalancer struggles with coarse moves involving LA County at $\alpha =
0.30$.
\ams{} converges to $\alpha^* \approx 0.22$ and achieves ACF 0.69 vs.\ 0.74
for fixed-$\alpha$ \ms{} --- a meaningful improvement.

\paragraph{Runtime.}
\ams{} runtime is within 2\% of fixed-$\alpha$ \ms{} runtime for all states.
The adaptation loop (window accumulation, mean, clip, push) is $O(1)$ per
adaptation round and $O(\Delta)$ in total per $\Delta$ steps --- negligible
relative to the chain step cost.

\subsection{Comparison to Optimal Fixed Alpha}

For each state, we ran a grid of fixed-$\alpha$ \ms{} chains at $\alpha \in
\{0.05, 0.10, 0.15, 0.20, 0.25, 0.30, 0.35, 0.40, 0.45, 0.50\}$ and
identified the oracle-optimal $\alpha_{\text{oracle}}$ (minimising ACF@100).

\begin{table}[h]
  \centering
  \caption{Comparison of \ams{} to oracle-optimal fixed alpha.}
  \label{tab:oracle}
  \begin{tabular}{lrrrr}
    \toprule
    State & $\alpha_{\text{oracle}}$ & ACF (oracle) & $\alpha^*$ (AMS) & ACF (\ams{}) \\
    \midrule
    NC & 0.25 & 0.49 & 0.28 & 0.50 \\
    TX & 0.35 & 0.57 & 0.33 & 0.58 \\
    IA & 0.15 & 0.40 & 0.17 & 0.41 \\
    WI & 0.25 & 0.52 & 0.26 & 0.53 \\
    CA & 0.20 & 0.68 & 0.22 & 0.69 \\
    \bottomrule
  \end{tabular}
\end{table}

\ams{} achieves ACF within 0.01--0.02 of the oracle-optimal fixed-$\alpha$
chain in all states.
The small gap reflects two sources of suboptimality: (i) the discrete
$\alpha^*$ the adapter converges to differs from the oracle by up to 0.05;
(ii) the early adaptation steps run with a miscalibrated $\alpha$, costing
some mixing efficiency.
These gaps are practically negligible.

\paragraph{Conclusion.}
\ams{} achieves near-oracle performance without any manual tuning.
It is a drop-in replacement for \texttt{--search multiscale}: any state that
benefits from multi-scale mixing benefits equally from adaptive multi-scale,
with no per-state calibration.
